% intro.tex (v6): four paragraphs, about 650 words. Citations use keys from refs_*.bib.
The active layer is the ground above permafrost that thaws in summer and refreezes in winter. Its thickness (ALT) sets the depth of seasonal ground movement under infrastructure, the volume of soil in which organic carbon is exposed to decomposition and the hydrological connection between surface and ground\cite{Biskaborn2019,ShiklomanovNelson2002}. ALT responds to the thawing degree-days of each summer, to snow and to the thermal state of the soil, so it is a yearly quantity: the same site can differ by tens of centimetres between a cool and a warm summer\cite{Nelson1997,Klene2001}. Observations, however, are sparse and uneven. Probing and ground-penetrating radar campaigns in Alaska, the Lena Delta and north-western Canada measure thaw depth at points, in different years and in different months, and most sites are visited once\cite{Streletskiy2025,Boike2024}.

Several circumpolar ALT products now exist at kilometre resolution. Statistical products fit machine-learning models to site observations with climate, terrain and soil covariates\cite{Aalto2018,Ran2022a,Wei2026,Liu2024}; process products drive a thermal model with remote-sensing land-surface temperature and snow\cite{Westermann2024,YiKimball2019}. Two features of the statistical products limit their use as yearly maps. First, the climate covariates are multi-year means and the targets are site means over all available years, so the output is a single map without a year, and reported validation statistics mix years and sites. Second, each product is validated on its own sites with its own split, often a random split of observations that places other years of the same site in the training set, so accuracies cannot be compared across products or with a physical model\cite{Roberts2017,Ploton2020,MeyerPebesma2022}. Reviews of the field have asked for validation that separates sites from years and for an explicit time definition of inputs and outputs\cite{Wadoux2021,Linnenbrink2024}.

Physical thaw models give a different route. The Stefan solution relates thaw depth to the square root of accumulated thawing degree-days through a single edaphic coefficient\cite{Stefan1891,RomanovskyOsterkamp1997,Riseborough2008}; the Kudryavtsev formulation adds snow and vegetation terms\cite{Kudryavtsev1974}. These models carry the year explicitly and transfer to new regions, but they cannot represent the terrain and soil differences between cells within a region. Recent studies have combined the two: physics outputs as model inputs, physics-generated pseudo-labels, and physics anchors with a learned residual\cite{Willard2022,Pilyugina2025,Gautam2025}. Each study reports one combination on one region and one split, so it is not known which combination helps, where, and under what observation density and covariate information.

Here we define the analysis unit as a cell-year: a 1 km cell with an ALT observation in a given year, forced with the ERA5-Land degree-days of that year accumulated to the day of observation\cite{munozsabater2021_era5_land}. With \nncyAll{} cell-years from three regions (Fig.~1) we ask four questions. (1) Does year-specific forcing improve on climatological forcing? (2) How does the ranking of physical, learned and combined methods change from random to spatial-block, year, block-by-year and region hold-out splits? (3) Where, and under which covariate and climate conditions, does a residual learner improve on a recalibrated Stefan model? (4) How do the methods compare with six published products on the same cells and years, and do they follow year-to-year change at sites observed in several years? All methods share the same inputs, splits and scoring rules, differences are reported with block-bootstrap confidence intervals, and the resulting 1 km maps carry the year of their forcing and an area of applicability.
